\chapter{粒子物理}

\problem{9.1}{试说明由于违反什么守恒律，使下列反应不能发生：
\begin{enumerate}
\item $\mathrm n\to\mathrm p+\mathrm e^-$；
\item $\mathrm n\to\pi^++\mathrm e^-+\nu_e$；
\item $\mathrm n\to\mathrm p+\pi^-$；
\item $\mathrm n\to\mathrm p+\gamma$。
\end{enumerate}}
% layout
\answerarea

\problem{9.2}{试计算当 $\gamma$ 射线和静止电子相互作用而产生正、负电子对的阈能。}
% layout
\answerarea
\newpage

\problem{9.3}{判断下列各反应中哪些是禁戒的，若是允许的反应，请说明是属于哪种相互作用。
\begin{enumerate}
\item $\mathrm p\to\pi^++\mathrm e^++\mathrm e^-$；
\item $\Lambda^0\to\mathrm p+\pi^-$；
\item $\mu^-\to\mathrm e^-+\bar\nu_e+\nu_\mu$；
\item $\mathrm n+\mathrm p\to\Sigma^++\Lambda^0$；
\item $\Sigma^+\to\mathrm p+\gamma$。
\end{enumerate}}
% layout
\answerarea

\problem{9.4}{讨论 $\Omega$ 超子可能的衰变方式，并说明为什么它只能通过弱作用衰变。}
% layout
\answerarea
\newpage

\problem{9.5}{用高能质子轰击液氢靶时，求产生 $\pi$ 介子的阈能。（氘核的质量为 $1875.6\ \mathrm{MeV}$。）}
% layout
\answerarea

\problem{9.6}{若规定 $\mathrm K^+$ 介子的奇异数为 $+1$，试定出 $\mathrm K^-$ 和 $\Sigma^+$ 粒子的奇异数。}
% layout
\answerarea
\newpage

\problem{9.7}{已知带电 $\pi$ 介子的寿命为 $2.6\times10^{-8}\ \mathrm s$，在实验室中测得飞行中 $\pi$ 介子的寿命为 $75\ \mathrm{ns}$，求 $\pi$ 介子的速度、总能量、动能和动量。}
% layout
\answerarea

\problem{9.8}{求静止 $\pi$ 介子衰变产生的 $\mu$ 子和中微子的动能。}
% layout
\answerarea
\newpage

\problem{9.9}{设一个介子衰变成两个自旋为零的中性 $\pi$ 介子，这个介子的自旋可能取哪些值？}
% layout
\answerarea

\problem{9.10}{一个高速电子与静止电子碰撞，求反应产生正电子的阈动量。}
% layout
\answerarea
\newpage

\problem{9.11}{用 $\pi^+$ 介子轰击静止质子时，测到有质量为 $1690\ \mathrm{MeV}$ 的共振态，求入射 $\pi^+$ 介子的最低动能。}
% layout
\answerarea

\problem{9.12}{$\pi^-$ 介子（$J^P=0^-$）和氘核（$J^P=1^+$）组成原子基态。$\pi^-$ 最终被氘核俘获，并发生反应 $\pi^-+\mathrm d\to\mathrm n+\mathrm n$，求：
\begin{enumerate}
\item 产生的中子对的轨道角动量；
\item 中子系统的总自旋角动量。
\end{enumerate}}
% layout
\answerarea
\newpage

\problem{9.13}{实验中观察到下列过程：$\Xi^-\to\Lambda^0+\pi^-$，试分析其同位旋、奇异数守恒情况，并说明这一过程是属哪种相互作用过程。}
% layout
\answerarea

\problem{9.14}{请写出强子 $\overline{\mathrm K}(\mathrm s\bar{\mathrm d})$、$\varphi(\mathrm s\bar{\mathrm s})$、$\Delta(\mathrm u\mathrm u\mathrm u)$、$\Xi(\mathrm s\mathrm s\mathrm d)$、$\overline\Omega(\bar{\mathrm s}\bar{\mathrm s}\bar{\mathrm s})$ 的电荷 $Q$、同位旋 $I$ 及 $I_3$、奇异数 $S$ 和重子数 $B$。}
% layout
\answerarea
\newpage

\problem{9.15}{阐述电子偶素 $^3\mathrm S_1$ 状态不可能衰变为两个光子的理由。}
% layout
\answerarea

\problem{9.16}{试证明由粲夸克和反粲夸克 $c\bar c$ 组成的介子不可能具有 $J^{PC}=1^{-+}$ 的状态。}
% layout
\answerarea
\newpage

\problem{9.17}{在光子、中微子、中子和电子中，试问：
\begin{enumerate}
\item 哪些不参与电磁相互作用？
\item 哪些不参与强相互作用？
\item 哪些不参与弱相互作用？
\end{enumerate}}
% layout
\answerarea

\problem{9.18}{在北京正负电子对撞机上的实验中，利用反应 $\mathrm e^++\mathrm e^-\to\mathrm J/\psi$，测到质量为 $3.1\ \mathrm{GeV}$ 的 $\mathrm J/\psi$ 粒子，求这一反应的阈能。如果以 $\mathrm e^+$ 撞击静止的 $\mathrm e^-$，则产生 $\mathrm J/\psi$ 的阈能为多大？}
% layout
\answerarea
\newpage

\problem{9.19}{$\mathrm J/\psi$ 粒子是由粲夸克和反粲夸克组成的，已知它的 $J^{PC}=1^{--}$，请给出这个粲夸克素的轨道角动量、自旋角动量、同位旋量子数。}
% layout
\answerarea

\problem{9.20}{已知 $\Xi^-$ 的 $Q=-1$，$B=1$，$S=-2$。请写出它的夸克组成，以及它的同位旋多重态中其他粒子的 $Q$、$B$ 和 $S$。}
% layout
\answerarea
\newpage

\problem{9.21}{$\rho^-$ 介子是由 $\pi^-$ 和 p 散射得到的介子共振态，其反应为 $\pi^-+\mathrm p\to\rho^-+\mathrm p$，试由守恒关系给出 $\rho^-$ 介子的同位旋、同位旋第三分量、奇异数和重子数。}
% layout
\answerarea

\problem{9.22}{实验测得 $\rho$ 介子的质量宽度为 $150\ \mathrm{MeV}$，求它的平均寿命。}
% layout
\answerarea
\newpage

\problem{9.23}{指出并说明下列哪些反应或衰变是守恒定律允许的：
\begin{enumerate}
\item $\mathrm p\to\pi^++\pi^0$；
\item $\pi^-+\mathrm p\to\mathrm K^++\Sigma^-$；
\item $\pi^-\to\mathrm e^-+\gamma$；
\item $\mathrm p\to\Sigma^++\pi^0$；
\item $\mathrm p+\mathrm p\to\mathrm n+\mathrm p+\pi^++\pi^-+\pi^0$。
\end{enumerate}}
% layout
\answerarea
\newpage

\problem{9.24}{衰变过程 $\Lambda\to\mathrm p+\pi^-$ 是一个弱作用过程，请给出判断的理由。}
% layout
\answerarea

\newpage

\problem{9.25}{试判断下列衰变过程哪些是可能发生的，哪些是不可能发生的，并加以说明。
\begin{enumerate}
\item $\pi^+\to\mu^++\nu_\mu$；
\item $\pi^-\to\pi^++\mathrm e^-+\mathrm e^-$；
\item $\pi^0\to\mathrm e^-+\mathrm p$；
\item $\mathrm p\to\mathrm n+\mathrm e^++\nu_e$；
\item $\mathrm n\to\mathrm p+\mathrm e^-+\bar\nu_e$。
\end{enumerate}}
% layout
\answerarea
\newpage

\problem{9.26}{根据夸克模型，说明强子中含有的价夸克或反夸克数。}
% layout
\answerarea

\newpage

\problem{9.27}{指出下列过程产生的 $\nu$ 中哪些是中微子，哪些是反中微子，并写出是属于哪种类型的：
\begin{enumerate}
\item $\Lambda^0\to\mathrm p+\mu^-+\nu$；
\item $\Sigma^+\to\mu^++\nu+\mathrm n$；
\item $\tau^+\to\mathrm e^++\nu+\nu$；
\item $\mu^-\to\mathrm e^-+\nu+\nu$；
\item $\mathrm K^+\to\mu^++\nu$。
\end{enumerate}}
% layout
\answerarea
\newpage

\problem{9.28}{假设某种玻色子的下列量子数都不为 0，试问这种玻色子和它的反粒子的哪一个量子数是相同的？
\begin{enumerate}
\item 奇异数；
\item 粲数；
\item 内禀宇称；
\item 磁矩；
\item 电荷。
\end{enumerate}}
% layout
\answerarea

\problem{9.29}{$\mathrm K^+$ 介子可以衰变为三个 $\pi$ 介子，$\mathrm K^+\to\pi^++\pi^-+\pi^+$，试计算 $\pi$ 介子的最大能量和平均能量。}
% layout
\answerarea
\newpage

\problem{9.30}{下列四组粒子中，哪一组是属于完备的同位旋多重态？
\begin{enumerate}
\item $(\rho^+,\rho^-)$；
\item $(\mathrm K^+,\mathrm K^-)$；
\item $(\mathrm K^0,\overline{\mathrm K}{}^0)$；
\item $(\Sigma^+,\Sigma^0,\Sigma^-)$。
\end{enumerate}}
% layout
\answerarea

\newpage

\problem{9.31}{费米子和反费米子系统的空间宇称 $\eta_P$ 与电荷宇称 $\eta_C$ 由系统的相对运动轨道角动量 $L$ 和总自旋 $S$ 决定，下列表达式中哪个是正确的？
\begin{enumerate}
\item $\eta_P=(-1)^L$，$\eta_C=(-1)^{L+S}$；
\item $\eta_P=(-1)^{L+1}$，$\eta_C=(-1)^{S+1}$；
\item $\eta_P=(-1)^{L+S}$，$\eta_C=(-1)^{L+1}$；
\item $\eta_P=(-1)^{L+1}$，$\eta_C=(-1)^{L+S}$。
\end{enumerate}}
% layout
\answerarea
\newpage

\problem{9.32}{已知共振态 $\psi(3686)$ 的夸克组成为 $2^3\mathrm S_1\ (c\bar c)$，其衰变总宽度 $\Gamma=277\ \mathrm{keV}$，衰变到正负电子对的分宽度 $\Gamma_{ee}=2.14\ \mathrm{keV}$。
\begin{enumerate}
\item 画出衰变末态为正负电子不变质量谱的形状（注意横坐标的比例），在不变质量谱上标出 $\psi(3686)$ 的质量、不变质量谱的半高全宽（FWHM）应该取何值（假定谱仪的质量分辨小于 $1\ \mathrm{keV}$）？
\item 计算该共振态的平均寿命。
\item 写出 $\psi(3686)$ 的 $I^G(J^{PC})$。
\end{enumerate}}
% layout
\answerarea
\newpage

\problem{9.33}{设能量在 GeV 的中微子的反应截面约为 $10^{-38}\ \mathrm{cm^2}$，试估算中微子在铀中的平均自由程。}
% layout
\answerarea

\problem{9.34}{由粒子表得到 $\Sigma$ 粒子的质量分别为 $1189.4\ \mathrm{MeV}$（$\Sigma^+$）、$1192.6\ \mathrm{MeV}$（$\Sigma^0$）和 $1197.4\ \mathrm{MeV}$（$\Sigma^-$），试由它们的夸克组成来估算 d 夸克和 u 夸克的质量差。}
% layout
\answerarea
\newpage

